% Gauss, Werke X/1 (1917), native93–97, printed86–90.
\SourceLeaf{093}\NoSourceHead
\SourceGroup{III}
\HistoricTitle{[ÜBER TERNÄRE QUADRATISCHE FORMEN.]}
\HistoricCitation{[Eintragung im \textsc{Leiste}, S. 111.]}
\NumberedItem{1}
\noindent\hspace*{\parindent}Zur Möglichkeit des Ausdrucks
\[
a\square+b=c\square
\]
\noindent wird erfordert:
\[
\begin{array}{rl}
1)&\text{dass }+bc\text{ ein Rest v. }a\text{ sei,}\\
2)&\text{dass }-ab\text{ ein Rest v. }c\text{ sei,}\\
3)&\text{dass }+ac\text{ ein Rest v. }b\text{ sei.}
\end{array}
\]
\noindent Allgemeiner, zur Möglichkeit der Gleichung
\[
a\square+b\square=c\square
\]
\noindent ist nöthig, dass
\[
\left.\begin{array}{rl}
1)&+ac\text{ Rest von }b\\
2)&-ab\text{ Rest von }c\\
3)&+bc\text{ Rest von }a
\end{array}\right\}\text{ sei.}
\]
\NumberedItem{2}
\noindent\hspace*{\parindent}Generalissime ad possibilitatem aequationis
\[
axx+byy+czz=0
\]
\noindent requiritur,
\clearpage
\SourceLeaf{094}\ArithmeticHead{87}{R}{ÜBER TERNÄRE QUADRATISCHE FORMEN.}
\noindent\hspace*{\parindent}I) ut non omnes $a,b,c$ idem signum habeant,\par
\noindent\hspace*{\parindent}II) ut
\[
\begin{array}{cccc}
&&&\shortstack{\scriptsize div. c.\\\scriptsize max.}\\
\multicolumn{1}{c|}{-bc}&\multicolumn{1}{c|}{\text{residua}}&\multicolumn{1}{c|}{a\times}&b,c\\
\multicolumn{1}{c|}{-ac}&\multicolumn{1}{c|}{\text{sint}}&\multicolumn{1}{c|}{b\times}&c,a\\
\multicolumn{1}{c|}{-ab}&\multicolumn{1}{c|}{\text{ipsorum}}&\multicolumn{1}{c|}{c\times}&a,b\\
\multicolumn{1}{c|}{\shortstack{\scriptsize a factt. suis\\\scriptsize quadratis\\\scriptsize liberati}}&\multicolumn{1}{c|}{}&\multicolumn{1}{c|}{}&
\end{array}
\]
\ClosingRule
\HistoricCitation{[Aus Scheda Ae, Varia, Julius 1800, S. 3—5.]}
\NumberedItem{3}
\begin{center}PROBLEMA.\\[6pt]
\textbf{Formam ternariam determinantis 0 in binariam transmutare.}\\[6pt]
\textit{Sol[utio] Casus I.}\end{center}
\noindent\hspace*{\parindent}Si omnes coefficientes formae adiunctae sunt 0.\par
\noindent\hspace*{\parindent}Sit proposita
\[
=\left(\begin{array}{ccc}a,&a',&a''\\b,&b',&b''\end{array}\right)
\]
\noindent habebuntque $a,a',a''$ idem signum. Sit $m$ ipsorum div. comm. max. eodem
signo acceptus, metieturque etiam ipsos $b,b',b''$. Tunc erit
\[
\begin{gathered}
axx+a'yy+a''zz+2byz+2b'xz+2b''xy\\
=m(\mathfrak{A}x+\mathfrak{B}y+\mathfrak{C}z)^2
\end{gathered}
\]
\noindent ipsique $\mathfrak{A},\mathfrak{B},\mathfrak{C}$ divisorem communem non habebunt.
Accipiantur
\[
\alpha,\beta,\gamma,\quad\alpha',\beta',\gamma',\quad\alpha'',\beta'',\gamma''
\]
\noindent ita ut sit
\[
\begin{aligned}
\alpha\mathfrak{A}+\beta\mathfrak{B}+\gamma\mathfrak{C}&=1,\\
\beta'\gamma''-\beta''\gamma'&=\mathfrak{A},\\
\gamma'\alpha''-\gamma''\alpha'&=\mathfrak{B},\\
\alpha'\beta''-\alpha''\beta'&=\mathfrak{C}
\end{aligned}
\]
\clearpage
\SourceLeaf{095}\ArithmeticHead{88}{L}{ARITHMETIK. NACHLASS.}
\noindent transibitque $f$ per substitutionem
\[
\begin{array}{ccc}\alpha,&\alpha',&\alpha''\\\beta,&\beta',&\beta''\\\gamma,&\gamma',&\gamma''\end{array}
\]
\noindent in formam
\[
\left(\begin{array}{ccc}m,&0,&0\\0,&0,&0\end{array}\right)
\]
\begin{center}\textit{Casus II.}\end{center}
\noindent\hspace*{\parindent}Si conditio in I apud formam propositam non habet locum, certo in forma
ipsi adiuncta locum habebit. Determinentur itaque $\alpha,\alpha',\alpha''$ etc. ita ut haec
adiuncta per substitutionem
\[
\begin{array}{ccc}\alpha,&\alpha',&\alpha''\\\beta,&\beta',&\beta''\\\gamma,&\gamma',&\gamma''\end{array}
\]
\noindent transeat in
\[
\left(\begin{array}{ccc}M,&0,&0\\0,&0,&0\end{array}\right)
\]
\noindent transibitque proposita per substitutionem illi adiunctam in formam talem
\[
\left(\begin{array}{ccc}0,&a',&a''\\b,&0,&0\end{array}\right)
\]
\noindent atque formae binariae $(a',b,a'')$ determinans erit $M$.
\NumberedItem{4}
\begin{center}PROBLEMA.\par\medskip\textbf{Solvere aequationem}
\[
axx+a'x'x'+a''x''x''+2(bx'x''+b'xx''+b''xx')=0,
\]
\textbf{determinans huius formae ternariae $f$ supponitur $=D$.}\end{center}
\noindent\hspace*{\parindent}Indoles solutionis consistit in inventione formae ternariae determinantis
$D^3$ cuius coefficientes omnes sint divisibiles per $D$ et quae sub $f$ contenta sit.
\clearpage
\SourceLeaf{096}\ArithmeticHead{89}{R}{ÜBER TERNÄRE QUADRATISCHE FORMEN.}
\noindent\hspace*{\parindent}1) Si aequatio est solubilis, talis forma dabitur.\par
\noindent Tunc formae propositae aequivalens inveniri potest forma talis
\[
\left(\begin{array}{ccc}0,&a',&a''\\b,&b',&0\end{array}\right)
\]
\noindent quae per substitutionem
\[
\begin{array}{ccc}1,&0,&0\\0,&b',&0\\0,&0,&a'b'\end{array}
\]
\noindent transit in
\[
D\left(\begin{array}{ccc}0,&1,&a'a''\\b,&1,&0\end{array}\right)
\]
\noindent proprietate requisita praeditam.
\NumberedItem{5}
\begin{center}\textbf{Algorithmus per quem si fieri potest ad talem formam pervenitur.}\end{center}
\noindent\hspace*{\parindent}I.\quad Proposita propria
\[
f=\left(\begin{array}{ccc}a,&a',&a''\\b,&b',&b''\end{array}\right)
\quad\text{det. }D=Mmm,
\]
\noindent adiuncta
\[
F=\left(\begin{array}{ccc}Am,&A'm,&A''m\\Bm,&B'm,&B''m\end{array}\right)
\quad\text{det. }=MMm^4.
\]
\noindent Tunc erit
\[
f\equiv(h,h',h'')^2\quad(\mathrm{mod.}\ m);
\]
\noindent fiat
\[
\begin{gathered}
\beta'\gamma''-\beta''\gamma'=h\text{ etc.}\\
ah+a'h'+a''h''=1,
\end{gathered}
\]
\noindent $f$ transit per substitutionem
\[
\begin{array}{ccc}m\alpha,&\beta,&\gamma\\m\alpha',&\beta',&\gamma'\\m\alpha'',&\beta'',&\gamma''\end{array}
\]
\noindent in formam talem
\par\smallskip\noindent{\scriptsize X1.}\hfill{\small 12}\par
\clearpage
\SourceLeaf{097}\ArithmeticHead{90}{L}{ARITHMETIK. NACHLASS.}
\[
\left(\begin{array}{ccc}mma,&a'm,&a''m\\bm,&b'm^2,&b''m^2\end{array}\right)
\]
\noindent eritque determinans formae
\[
\left(\begin{array}{ccc}ma,&a',&a''\\b,&mb',&mb''\end{array}\right)
\]
\noindent $Mm$. Conditio possibilitatis aequationis
\[
\begin{aligned}
f&=\left(\begin{array}{ccc}a,&a',&a''\\b,&b',&b''\end{array}\right)
\mathinner{.\,.\,.}\,Mmm\,\mathinner{.\,.\,.}
\left(\begin{array}{ccc}Am,&A'm,&A''m\\Bm,&B'm,&B''m\end{array}\right)\\[6pt]
F&=\left(\begin{array}{ccc}A,&A',&A''\\B,&B',&B''\end{array}\right)
\mathinner{.\,.\,.}\,MMm\,\mathinner{.\,.\,.}
\left(\begin{array}{ccc}aM,&a'M,&a''M\\bM,&b'M,&b''M\end{array}\right)
\end{aligned}
\]
\[
\boxed{\begin{aligned}
\frac{Mf}{\square}\;&R\;\frac{m}{\square}\\[4pt]
\frac{mF}{\square}\;&R\;\frac{M}{\square}\\[4pt]
-\frac{MmfF}{\square}\;&R\;\text{div. comm. max. ipsorum }\frac{m}{\square},\frac{M}{\square}
\end{aligned}}
\]
\begin{center}\small BEMERKUNGEN.\end{center}
\begingroup\small
\noindent\hspace*{\parindent}Die Aufzeichnung [1.] stammt aus sehr früher Zeit, etwa 1796, die Aufzeichnung [2.] ist in anderer
Schrift, offenbar später, hinzugeschrieben. Zu [3.] vergleiche man die artt. 281 (Schluß) und 299 der
\textit{Disquisitiones Arithmeticae} [Werke I, S. 323 und 358]. In [4.] werden die Bedingungen für die Auflösbarkeit
einer ternären quadratischen Gleichung gegeben. Diese sind von H. St. \textsc{Smith} (Proc. of the Royal
Society London \textbf{13}, 1864, S. 110, Papers I, S. 410) aufgestellt und von \textsc{Arnold Meyer} (\textsc{Crelles} Journal
f. Mathematik \textbf{98}, 1885, S. 177) bewiesen worden, und stimmen mit den in der vorstehenden Aufzeichnung
von \textsc{Gauss} enthaltenen vollkommen überein. Vergl. den Artikel \textbf{12} des \textsc{Bachmann}schen Aufsatzes »Über
\textsc{Gauss}’ zahlentheoretische Arbeiten«. Im vorstehenden bedeuten $\frac{m}{\square}$ u.\,s.\,w. die von ihrem größten quadratischen
Teiler befreiten Zahlen.
\par\hfill\textsc{Schlesinger.}\par
\endgroup
\ClosingRule
\clearpage
